\documentclass[10pt,a4paper]{amsart}
\usepackage[margin=19mm]{geometry}
\usepackage{amsmath,amssymb,mathtools}
\usepackage[scr=rsfs,cal=euler]{mathalfa}
\usepackage{xfrac}
\usepackage[unicode,hidelinks,bookmarks=false]{hyperref}
\newcommand{\ie}{i.e.\ }
\newcommand{\cf}{cf.\ }
\newcommand{\resp}{respectively\ }
\newcommand{\Subsection}{\subsection}
\providecommand{\nobreakdash}{\nobreak\hbox{-}}
\setlength{\emergencystretch}{2em}
\multlinegap0pt
\usepackage{amssymb}
\usepackage[shortlabels]{enumitem}
\usepackage{csquotes}
\usepackage{tikz}
\usepackage{xcolor}
\newtheorem{prop}{Proposition}
\newtheorem{lemma}{Lemma}
\newcommand{\Id}{\mathsf{Id}}
\DeclareMathOperator{\Lip}{Lip}
\newcommand{\MM}{\mathcal{M}}
\renewcommand{\SS}{\mathscr{S}}
\newcommand{\cZ}{\mathcal{Z}}
\newcommand{\ee}{\mathsf{e}}
\renewcommand{\ge}{\geqslant}
\renewcommand{\le}{\leqslant}
\newcommand{\sub}{\mathscr{C}}
\renewcommand{\subset}{\subseteq}
\title{De-H\"oldering factorization}
\author{Assaf Naor}
\date{Source: arXiv:2609.07564v2 (CC BY 4.0)}
\begin{document}
\maketitle

\noindent\emph{Source and licence.} De-H\"oldering factorization. Version arXiv:2609.07564v2 (CC BY 4.0), distributed by arXiv under the Creative Commons Attribution 4.0 International licence, http://creativecommons.org/licenses/by/4.0/, as recorded in the arXiv licence metadata for this exact version and shown on the abstract page https://arxiv.org/abs/2609.07564v2. Full licence notice: \texttt{licenses/arxiv-2609.07564-CC-BY-4.0.txt}.\par\smallskip

\noindent\textit{Editorial scope.} Counted unit: the article's prop ext (source0.tex lines 333--339). The article's complete original proof of it is delivered with it (source0.tex lines 571--620, one \texttt{proof} environment of 50 lines, which the article places away from the statement). No mathematics is written, repaired or re-derived: the statement and the proof are the article's own lines, in the article's own notation, and no retained line is substituted or re-flowed.  Retained uncounted as the source-local context the proof needs: lines 551--560.  Nothing was omitted from any retained span, so the package carries no omission record.  No retained line contains a citation, so the wrapper carries no bibliography and no bibliography entry.  No retained line contains a figure environment, an image-inclusion command or drawing code, so no image is delivered and the package declares no asset dependency.  Editorial formatting only, in addition: an AMS \texttt{amsart} preamble that restates the article's own declarations and macros used by the retained lines, namely the environments \texttt{prop}, \texttt{lemma} on separate counters, together with the standard packages those lines need.  The retained environments print in this document as Proposition 1, Lemma 1 in order of appearance, whereas the article prints the same items with the numbering of its own full document.  The complete native source of the article as distributed in the arXiv e-print for this exact version is bundled unchanged as \texttt{sources/209640/source0.tex}.
\begin{lemma}\label{lem:two scales} Let $\omega:(0,\infty)\to (0,\infty)$ be nondecreasing and concave, with $\lim_{t\to 0^+}\omega(t)=0$. Fix  metric spaces $(\MM,d_\MM), (\cZ,d_\cZ)$ such that $\cZ$ has a conical geodesic bicombing. If $\sub\subset \MM$ is closed and $f:\sub\to \cZ$ satisfies:
\begin{equation}\label{eq:omega modulus}
\forall x,y\in \sub,\qquad d_\cZ\big(f(x),f(y)\big)\le \omega \big(d_\MM(x,y)\big),
\end{equation}
then there exists $F:\MM\to \cZ$ such that $F(a)=f(a)$ for every $a\in \sub$,  and for every distinct $x,y\in \MM$ we have:
\begin{equation}\label{eq:omega extended}
d_\cZ\big(F(x),F(y)\big)\lesssim \ee(\MM;\cZ)\frac{\omega\left(d_\MM(x,y)+d_\MM(\{x,y\},\sub)\right)}{d_\MM(x,y)+d_\MM(\{x,y\},\sub)}d_\MM(x,y).
\end{equation}
Furthermore, if $\cZ$ is a Banach space, then $F$ can be taken to depend linearly on $f$, with $\ee(\MM;\cZ)$ replaced by $\ee_l(\MM;\cZ)$ in~\eqref{eq:omega extended}.  
\end{lemma}

\begin{prop}\label{prop:def to ext} Let $(\MM,d_\MM), (\cZ,d_\cZ)$ be metric spaces such that $\cZ$ has a conical geodesic bicombing. If the identity mapping  on $\MM$  admits a de-H\"oldering factorization through $\cZ$, then   $S(\cZ)\subset S(\MM)$. 

Quantitatively, if the identity mapping on $\MM$ admits  a $K$-de-$\theta$-H\"oldering factorization through  $\cZ$ for some $0<\theta\le 1$ and $K\ge 1$, then $\ee(\SS,\MM)^\theta\lesssim K\ee(\SS,\cZ)$ for every metric space $(\SS,d_\SS)$.\footnote{We use the following  asymptotic notation, in addition to  the usual $O(\cdot)$ notation. Given $a,b>0$, by writing
$a\lesssim b$ or $b\gtrsim a$ we mean that $a\le \kappa b$ for some
universal constant $\kappa>0$, and $a\asymp b$
stands for $(a\lesssim b) \wedge  (b\lesssim a)$. }
\end{prop}

\begin{proof}[Deduction of Proposition~\ref{prop:def to ext}  from Lemma~\ref{lem:two scales}] Fix $0<\theta\le 1$ and $K\ge 1$. Suppose that  $\Id_\MM:\MM\to \MM$ admits a $K$-de-$\theta$-H\"oldering factorization through  $\cZ$. This means that there are $g:\MM\to \cZ$ and $h:\cZ\to \MM$ with: 
\begin{equation}\label{eq:composition identity}
h\circ g=\Id_\MM,
\end{equation} 
such that the following two conditions hold for some $\alpha>0$: 
\begin{equation}\label{eq:factor}
\left\{\begin{array}{ll}
\forall x,y\in \MM,\qquad d_\cZ\big(g(x),g(y)\big)\le \alpha d_\MM(x,y)^{\theta},\\ \forall z,w\in \cZ,\qquad d_\MM\big(h(z),h(w)\big)^\theta \le \frac{K}{\alpha}  \max \big\{d_\cZ(z,w),  d_\cZ\big(\{z,w\},g(\MM)\big)^{1-\theta}d_\cZ(z,w)^\theta\big\}.\end{array}\right. 
\end{equation}
 Letting $(\SS,d_\SS)$ be a metric space satisfying $\ee(\SS;\cZ)<\infty$, our goal is to deduce that:   
 \begin{equation}\label{eq:restate extension goal}
 \ee(\SS;\MM)^\theta\lesssim K\ee(\SS;\cZ).
 \end{equation} 

Suppose that  $\emptyset\neq \sub\subset \SS$ is closed and that $f:\sub\to \MM$ is Lipschitz. By the first part of~\eqref{eq:factor} we have:
$$
\forall s,t\in \sub,\qquad d_\cZ\big(g\circ f(s),g\circ f(t)\big)\le \alpha \|f\|_{\Lip(\sub;\MM)}^\theta d_\SS(s,t)^\theta. 
$$
Consequently,  Lemma~\ref{lem:two scales} applied to $g\circ f:\sub\to \cZ$, provides $\Gamma:\SS\to \cZ$  that satisfies:
\begin{equation}\label{eq:Gamma composition}
\forall s\in \sub,\qquad \Gamma(s)=g\circ f(s),
\end{equation}
as well as: 
\begin{equation}\label{eq:Gamma}
\forall s,t\in \SS,\qquad  d_\cZ\big(\Gamma(s),\Gamma(t)\big)\lesssim   \frac{\ee(\SS;\cZ)\alpha \|f\|_{\Lip(\sub;\MM)}^\theta d_\SS(s,t) }{\left(d_\SS(s,t)^{\phantom{2}}\!\!\!\!+d_\SS(\{s,t\},\sub)\right)^{1-\theta}}\le \ee(\SS;\cZ)\alpha \|f\|_{\Lip(\sub;\MM)}^\theta d_\SS(s,t)^\theta.
\end{equation}

 Observe that:
$$
\forall s\in \SS,\qquad d_\cZ\big(\Gamma(s),g(\MM)\big)\le  d_\cZ\big(\Gamma(s),g\circ f(\sub)\big)\stackrel{\eqref{eq:Gamma composition}}{=} d_\cZ\big(\Gamma(s),\Gamma(\sub)\big)\stackrel{\eqref{eq:Gamma}}{\lesssim} \ee(\SS;\cZ) \alpha \|f\|_{\Lip(\sub;\MM)}^\theta  d_\SS(s,\sub)^\theta, 
$$
where the first step holds because $f(\sub)\subset \MM$. Consequently:
\begin{equation}\label{eq:st}
\forall s,t\in \SS,\qquad d_\cZ\big(\{\Gamma(s),\Gamma(t)\},g(\MM)\big)\lesssim  \ee(\SS;\cZ) \alpha \|f\|_{\Lip(\sub;\MM)}^\theta   d_\SS(\{s,t\},\sub)^\theta,
\end{equation}
and therefore the following estimate holds for every $s,t\in \SS$:
$$
\max\Big\{d_\cZ\big(\Gamma(s),\Gamma(t)\big),d_\cZ\big(\{\Gamma(s),\Gamma(t)\},g(\MM)\big)^{1-\theta}d_\cZ\big(\Gamma(s),\Gamma(t)\big)^\theta\Big\}
\stackrel{\eqref{eq:Gamma}\wedge\eqref{eq:st}}{ \lesssim} \ee(\SS;\cZ)\alpha \|f\|_{\Lip(\sub;\MM)}^\theta  d_\SS(s,t)^\theta.
$$
In combination with~\eqref{eq:factor}, this implies that: 
$$
\forall s,t\in \SS,\qquad d_\MM\big(h\circ \Gamma(s),h\circ \Gamma(t)\big)^\theta\lesssim K\ee(\SS;\cZ) L^\theta d_\SS(s,t)^\theta.
$$
In other words, the Lipschitz constant of $h\circ \Gamma:\SS\to \MM$ satisfies: 
$$
\|h\circ  \Gamma\|_{\Lip(\SS;\MM)}^\theta\lesssim K\ee(\SS;\cZ)\|f\|_{\Lip(\sub;\MM)}^\theta.
$$
So, $\ee(\SS,\sub;\MM)^\theta\lesssim K\ee(\SS;\cZ)$, as   $h\circ \Gamma$ extends $f$ thanks to~\eqref{eq:composition identity} and~\eqref{eq:Gamma composition},  and $f$ is an arbitrary Lipschitz function from $\sub$ to $\MM$. Because $\sub$ is an arbitrary closed subset of $\SS$, the desired estimate~\eqref{eq:restate extension goal} follows.
\end{proof}

\end{document}
